% ARCHIVED at v3.66 (2026-09-22, author's instruction): Chapter 6, subsubsection 6.1.2.6 "Endpoint and Per-Step Projection" (06_results.tex lines 571-677 of v3.65), including Fig 6.5 fig:projector-cost and Table 6.3 tab:hf-ladder
% Verbatim copy of the v3.65 text, line-for-line, for reference only. NOT built. Do not edit.

\subsubsection{Endpoint and Per-Step Projection}\label{sec:res:avoiding:projectors}

The models above were read twice, unprojected and under the per-step projector. This subsection holds the
model fixed instead and varies the projector. It is made on one generative model, and that is a limit of
the data rather than a choice: on this benchmark endpoint projection was run on the two average-velocity
models and never on instantaneous-velocity matching or on the diffusion baseline. At the budgets this
section is locked to, only the analytic objective has it over all five training seeds, so the comparison
below is made there. The cross-model version of it is on D3IL-aligning, where all three flow-based models
carry both projectors (\autoref{tab:va-projection-models}).
\guard{No endpoint-projection evaluation of instantaneous-velocity matching or of the diffusion baseline
exists on D3IL-avoiding, so no statement in this subsection extends to those two models.
Recorded in \texttt{data\_status/PENDING\_20260922\_all\_lacking\_runs.md}.}

Endpoint projection \parencite{li2025hardflow} and per-step projection were run on the same analytic average-velocity matching
model, with the same solver and constraints, in one evaluation. Endpoint projection solves more
programs per control step and still takes about half the time: $0.0745$ against $0.1478$\,s at $\nfe=3$,
and $0.1304$ against $0.2268$\,s at $\nfe=5$ (\autoref{fig:projector-cost}). Its predicted endpoint is
close to the feasible set at every step, so the solver needs few iterations.

\begin{figure}[htpb]
  \centering
  \includegraphics[width=\textwidth]{fig_avoiding_projector_cost}
  \caption[The cost of the two projection methods]{The time to compute one action
  (\autoref{sec:setup:metrics:avoiding}) under per-step and under endpoint projection, against the step
  budget; analytic average-velocity matching (MeanFM), temporal U-Net of $4.0$\,M parameters, activation
  threshold $1.0$, temporal consistency, tightened constraints, \textbf{four training seeds (7--10)
  $\times$ three geometries $\times$ two episodes}, 24 episodes per configuration. Both methods are read
  on the same four seeds, so this is a matched multi-seed comparison and not a single-seed one; seed 6 is
  the one absent from this evaluation. The shaded budget has a single guiding step; the hollow marker is
  the endpoint series there.}
  \label{fig:projector-cost}
\end{figure}

At $\nfe=3$, endpoint projection reaches a success rate with constraint satisfaction of $1.000$ with no
violations, in $58.88$ control steps at $0.0745$\,s per step. The \ac{DPCC} baseline reaches $1.000$ in
$70.13$ steps at $0.5534$\,s per step.
\dataref{temp/0609/I/batch\_avoiding\_combined\_20260906\_125724, job 25444; DA\_20260906 \S3}
\provisional{This comparison uses four training seeds --- 7, 8, 9 and 10 --- $\times$ three geometries
$\times$ two episodes, 24 episodes per configuration; both methods are read on the same four, and seed 6
is absent from the evaluation despite the run tag naming it. The difference in success with per-step
projection at the same budget, $1.000$ against $0.917$, is two episodes. The evaluation at five training
seeds and twenty episodes is pending.}

That result is at $\nfe=3$. The budget at which this task is actually run is lower, and there endpoint
projection does not execute its own arithmetic at all. With the activation threshold at $0.5$ the
guiding step count (\autoref{sec:res:constraints:degenerate}) is zero at $\nfe=1$ and at $\nfe=2$: the
only active step is the terminal one, where the predicted endpoint is the finished sample and the
method reduces to projection after sampling. This is measured, not inferred --- the evaluation records
one program solved per plan at both budgets, against three at $\nfe=5$ and five at $\nfe=10$, matching
the predicted count exactly. A run at $\nfe=2$ with both methods given the same number of candidate
plans produced bit-identical rollouts, at $2.45$ times the cost.

Across the budgets at which it does run, endpoint projection shows no consistent advantage on this task
(\autoref{tab:hf-ladder}). Untightened it reaches higher success with constraint satisfaction than all
three selection rules of per-step projection only at $\nfe=10$; on tightened constraints it does so only
at $\nfe=1$, where it runs no endpoint-projection arithmetic, and at $\nfe=10$ it is ahead of none of
them. Its one systematic advantage is control steps, and the mechanism is not the projection: with four
candidate plans the cumulative-projection-cost rule sometimes selects a stalled plan of 97 control steps,
and endpoint projection, given one candidate, cannot make that selection.

\begin{table}[htpb]
  \caption[Endpoint projection across the step budget]{Endpoint projection against per-step projection on
  obstacle avoidance, across the step budget. Success with constraint satisfaction, counted as how many of
  the three selection rules of per-step projection endpoint projection is ahead of. Analytic
  average-velocity matching, activation threshold $0.5$, \textbf{five training seeds (6--10) $\times$
  three geometries $\times$ two episodes}, so every entry rests on 30 episodes and 15 seed--geometry
  cells. Guiding steps are the predicted count, confirmed against the recorded number of programs solved
  per plan. The per-budget success, control-step and time numbers this table reduces to counts are in
  \autoref{tab:app:hf-ladder-detail}.}
  \label{tab:hf-ladder}
  \centering
  \small
  \begin{tabularx}{\linewidth}{@{}rrL L@{}}
    \toprule
    $\nfe$ & Guiding steps & Untightened & Tightened \\
    \midrule
    1  & 0 & ahead of 2 of 3 & ahead of 3 of 3 \\
    2  & 0 & ahead of 2 of 3 & ahead of 0 of 3, level with 1 \\
    5  & 2 & ahead of 2 of 3 & ahead of 2 of 3, level with 1 \\
    10 & 4 & ahead of 3 of 3 & ahead of 0 of 3, level with 2 \\
    \bottomrule
  \end{tabularx}
\end{table}

\dataref{temp/1108/Revised\_2/batch\_avoiding\_combined\_20260811\_221322; DA\_20260824 \S2--\S4}
\guard{On this ladder endpoint projection was run with one candidate plan and per-step projection with
four, so every row in which endpoint projection is ahead understates it and every row in which it is
behind is confounded. The time columns of that evaluation are unusable for the same reason. The
time comparison at $\nfe=3$ and $\nfe=5$ quoted at the start of this subsection gave both methods four
candidate plans.}

The consequence for this task is that endpoint projection is not needed here, and the reason is the
budget rather than the method. At one network evaluation the two average-velocity models already reach
$0.967$ and $1.000$ success with constraint satisfaction at about $18$\,ms per control step with the
per-step projector, and instantaneous-velocity matching reaches $1.000$ at one and at two
(\autoref{tab:avoiding-dpcc-protocol}); there is no shortfall left for a second projection method to
close. Endpoint projection needs sampling steps to guide, and at
one and two evaluations it has none or one (\autoref{sec:method:degenerate},
\autoref{sec:res:constraints:degenerate}), so at the budget this benchmark is run at it does not execute
its own arithmetic --- \autoref{tab:hf-ladder} prints those two rows for completeness only, and the
$\nfe=2$ run returned rollouts identical to per-step projection at $2.45$ times the cost. Making it
available means raising the budget to five or ten evaluations, and \autoref{fig:projector-cost} shows
what that costs: per control step the projector, not the network, is the expense on this task. Raising
the budget therefore gives up precisely the advantage that decides the benchmark, in exchange for a
method that shows no consistent gain once it does run. Endpoint projection is kept in the thesis, and
tested where it can act, on D3IL-aligning and the quadrotor scenes, which are run at larger budgets.

